\documentclass[12pt]{article}

\setlength{\parindent}{0pt}
\setcounter{secnumdepth}{3}
\usepackage[margin=1in]{geometry}

% packages
\usepackage{amsmath, amsfonts, amsthm, amssymb}
\usepackage{graphicx, float}
\usepackage{mathtools}
\usepackage{titlesec}
\usepackage{interval}
\usepackage{hyperref}
\usepackage{siunitx}
\usepackage{titling}
\usepackage{vwcol}
\usepackage{setspace}
\usepackage{empheq}
\usepackage{cancel}
\usepackage{esdiff}
\usepackage{multicol}
\usepackage{mdframed}
\usepackage{esdiff}
\usepackage{tikzsymbols}
\usepackage{multicol}
\usepackage{tikz}
\usepackage{varwidth}
\usepackage{pgfplots}
\usepackage{fancyhdr}
\usepackage{enumitem}
\usepackage{tcolorbox}

\intervalconfig {
	soft open fences
}

% header
\pagestyle{fancy}
\fancyhf{}
\lhead{Ethan Fan}
\rhead{AP Calculus BC}
\cfoot{\thepage}

% section format
\titleformat{\section}
  {\large \bfseries}
  {}
  {0em}
  {}
\titlespacing*{\section}{0pt}{0.5em}{0.3em}

\titleformat{\subsection}[runin]
  {\bfseries}
  {}
  {0em}
  {}
\titlespacing*{\subsection}{0pt}{0pt}{0em}

\titleformat{\subsubsection}[runin]
  {\itshape}
  {(\roman{subsubsection})}
  {0em}
  {}

% commands
\newcommand{\alignedintertext}[1]{%
  \noalign{%
    \vskip\belowdisplayshortskip
    \vtop{\hsize=\linewidth#1\par
    \expandafter}%
    \expandafter\prevdepth\the\prevdepth
  }%
}

\newcommand*{\paren}[1]{\ensuremath\left(#1\right)}
\newcommand*{\limit}[2][x]{\ensuremath{\displaystyle\lim_{#1\to#2}}}
\newcommand*{\Limit}[3][x]{\ensuremath{\displaystyle\lim_{#1\to#2}\left[#3\right]}}
\newcommand*{\deriv}[1][x]{\ensuremath{\dfrac{\mathrm{d}}{\mathrm{d}#1}}}
\newcommand*{\Deriv}[2][x]{\ensuremath{\dfrac{\mathrm{d}}{\mathrm{d}#1}\left[#2\right]}}
\newcommand{\ddx}{\dfrac{d}{dx}}
\newcommand{\dydx}{\dfrac{dy}{dx}}
\newcommand{\dydt}{\dfrac{dy}{dt}}
\newcommand{\dxdt}{\dfrac{dx}{dt}}
\newcommand*{\iinteg}[2][x]{\ensuremath{\displaystyle\int #2\;\mathrm{d}#1}}
\newcommand*{\dinteg}[4][x]{\ensuremath{\displaystyle\int_{#2}^{#3}#4\;\mathrm{d}#1}}
\newcommand*{\abs}[1]{\ensuremath{\left|#1\right|}}
\newcommand*{\eps}{\varepsilon}
\newcommand*{\real}{\ensuremath\mathbb{R}}
\newcommand*{\integer}{\ensuremath\mathbb{Z}}
\newcommand*{\posint}{\ensuremath\mathbb{Z^+}}
\newcommand*{\cbrt}[1]{\ensuremath{\sqrt[3]{#1}}}
\newcommand*{\lheqzero}{\ensuremath{\underset{\text{L'H}}{\overset{\left[\frac00\right]}{=}}}}
\newcommand*{\lheqinfty}{\ensuremath{\underset{\text{L'H}}{\overset{\left[\frac{\infty}{\infty}\right]}{=}}}}
\newcommand{\tor}{\text{ or }}
\newcommand{\floor}[1]{\lfloor #1 \rfloor}

\newcommand{\vem}{\vspace{0.5em}}
\newcommand{\example}[1]{\section{Example #1}}
\newcommand{\problem}[1]{\section{Problem #1}}
\newcommand{\subproblem}[1]{\subsection{(#1)} \,}

\DeclareMathOperator{\dne}{DNE}
\DeclareMathOperator{\sgn}{sgn}

\newcommand{\definition}[1]{\begin{tcolorbox}[colback=red!5!white,colframe=red!75!black,parbox=false] #1 \end{tcolorbox}}

\newcommand{\theorem}[2]{\begin{tcolorbox}[title={#1},colback=blue!5!white,colframe=blue!75!black,parbox=false] #2 \end{tcolorbox}}

\allowdisplaybreaks
% \postdisplaypenalty=100000

% document
\begin{document}

\vspace*{0cm}

% opening
\begin{center}
    {\Large \bfseries Problem Set 16} \\[0.5cm]
    {\large Triangle Substitution} \\[0.35cm]
    {\normalsize October 13, 2026}
\end{center}

% problems

\problem{6}

\subproblem{a}
\begin{gather*}
    \int \frac{x}{\sqrt{9-x^2}} \, dx =
    \begin{bmatrix}
        u=9-x^2\\
        du= -2x\, dx
    \end{bmatrix} -\frac{1}{2}\int \frac{du}{\sqrt{u}} \\
    = -\frac{1}{2}\cdot2\sqrt{u} = \boxed{-\sqrt{9-x^2}+C}
\end{gather*}

\subproblem{b}
\begin{gather*}
    \int \frac{x^3}{\sqrt{9-x^2}} \, dx  =
    \begin{bmatrix}
        u=9-x^2\\
        du= -2x\, dx
    \end{bmatrix} \frac{1}{2}\int \frac{u-9}{\sqrt{u}} \, du \\
    =\frac{1}{2} \int \sqrt u - \frac{9}{\sqrt{u}} \, du = -9\sqrt{u} + \frac{1}{3}u^\frac{3}{2} \\
    = \boxed{\frac{1}{3}(9-x^2)^\frac{3}{2} - 9\sqrt{9-x^2} +C}
\end{gather*}

\subproblem{c}
\begin{gather*}
    \int \sqrt{9-x^2} \, dx = \int \sqrt{9-(3\sin \theta )^2} 3\cos x\, d\theta \\
    = 9 \int \cos ^2 \theta \, d\theta  = \frac{9}{2} \int \cos 2 \theta +1 \, d \theta \\
    = \frac{9}{2} \sin \theta \cos \theta  +\frac{9\theta}{2} +C\\
    = \boxed{\frac{9}{2}\arcsin \frac{x}{3} + \frac{x\sqrt{9-x^2}}{2}+C}
\end{gather*}

\subproblem{d}
\begin{gather*}
    \int \frac{dx}{9+x^2} =
    \begin{bmatrix}
        u=\frac{x}{3}\\
        du=\frac{1}{3}\, dx
    \end{bmatrix} \frac{1}{3} \int \frac{1}{1+u^2} \,du \\
    =\frac{1}{3} \arctan u + C= \boxed{\frac{1}{3}\arctan \frac{x}{3}+C}
\end{gather*}

\problem{7}

\subproblem{b}
\begin{gather*}
    \int \frac{dx}{(x^2-4x+13)^\frac{3}{2}}
    = \int \frac{dx}{((x-2)^2+9)^\frac{3}{2}}\\
    =\begin{bmatrix}
        u=x-2\\
        du=dx
    \end{bmatrix}
    \int \frac{du}{(u^2+9)^\frac{3}{2}}\\
    = \int \frac{3\sec^2 \theta\, d\theta }{27\sec^3 \theta }
    = \frac{1}{9} \int \cos \theta \, d\theta \\
    =\frac{1}{9} \sin \theta + C = \frac{1}{9} \cdot \frac{u}{\sqrt{u^2+9}}+C\\
    =\boxed{ \frac{x-2}{9\sqrt{x^2-4x+13}} + C}
\end{gather*}

\problem{9}

\subproblem{a}
\begin{gather*}
    \sin \frac{\theta}{2}=\frac{t}{\sqrt{1+t^2}},\, \cos \frac{\theta}{2}= \frac{1}{\sqrt{1+t^2}}\\
    \sin \theta = 2 \cdot \frac{t}{\sqrt{1+t^2}} \cdot \frac{1}{\sqrt{1+t^2}} = \boxed{\frac{2t}{1+t^2}}\\
    \cos \theta = \paren{\frac{1}{\sqrt{1+t^2}}}^2 - \paren{\frac{t}{\sqrt{1+t^2}}}^2 = \boxed{\frac{1-t^2}{1+t^2}}\\
    \theta = 2 \arctan t \implies \boxed{d \theta = \frac{2\, dt}{1+t^2}}
\end{gather*}

\subproblem{b}
\begin{gather*}
    \int \frac{d\theta}{2+\cos \theta } 
    = \int \frac{2dt}{1+t^2}\cdot \frac{1+t^2}{3+t^2} = \int \frac{2}{t^2+3}\, dt\\
    = \frac{2}{\sqrt{3}} \arctan  \frac{t}{\sqrt{3}} + C =  \boxed{\frac{2}{\sqrt{3}} \arctan \paren{\frac{1}{\sqrt{3}} \tan \frac{\theta}{2}}  + C}
\end{gather*}

\subproblem{c}
It is impossible, as the integrand is always positive. The mistake is due to the domain restriction of the half-angle formulas. 
\begin{gather*}
    \int_0^{2\pi} \frac{d\theta}{2+\cos \theta } = \lim_{x \to \pi^-} F(x) - \lim_{x\to -\pi^+} F(x) = 2\cdot \frac{\pi}{\sqrt{3}} = \boxed{\frac{2\pi}{\sqrt{3}}}
\end{gather*}

\subproblem{d}
\begin{gather*}
    a>b \implies \frac{1}{a+b\cos \theta} >0\\
    \int_0^{2\pi} \frac{d\theta}{a+b\cos \theta } =  \int_{-\pi}^\pi \frac{d\theta}{a+b\cos \theta} \\
    = \int_{-\infty}^\infty \frac{2\, dt}{1+t^2} \cdot \frac{1+t^2}{(a+b)+(a-b)t^2}\\
    =\int_{-\infty}^\infty \frac{2\, dt}{(a+b)+(a-b)t^2} = \frac{2}{a-b}\int_{-\infty}^\infty \frac{dt}{\frac{a+b}{a-b}+t^2}\\
    = \begin{bmatrix}
        u=t\sqrt{\frac{a-b}{a+b}} \\
        du=dt\sqrt{\frac{a-b}{a+b}}
    \end{bmatrix}  \frac{2}{a-b} \sqrt{\frac{a-b}{a+b}} \left[ \arctan t\sqrt{\frac{a-b}{a+b}} \right]^\infty_{-\infty}\\
    =\boxed{\frac{2\pi}{\sqrt{a^2-b^2}}} 
\end{gather*}








\end{document}

